%This work is licensed under the Creative Commons
%Attribution-ShareAlike 4.0 International License. To view a copy of
%this license, visit http://creativecommons.org/licenses/by-sa/4.0/ or
%send a letter to Creative Commons, PO Box 1866, Mountain View, CA
%94042, USA.

%This work is licensed under the Creative Commons
%Attribution-ShareAlike 4.0 International License. To view a copy of
%this license, visit http://creativecommons.org/licenses/by-sa/4.0/ or
%send a letter to Creative Commons, PO Box 1866, Mountain View, CA
%94042, USA.

%\documentclass[gray,handout, pdftex, 11pt]{beamer}
%\documentclass[handout, pdftex, 11pt]{beamer}

\documentclass[pdftex, 11pt]{beamer}

\usepackage[utf8]{inputenc}
\usepackage[T1]{fontenc}
\usepackage{lmodern}
%\usepackage[italian]{babel}
\usepackage{graphicx}
\usepackage{listings}
\usepackage{microtype}
\usepackage{acronym}
\usepackage{array}
\usepackage{tikz}
\usetikzlibrary{shapes, chains, scopes, shadows, positioning, arrows,
  decorations.pathmorphing, calc}

\colorlet{c1}{green!20}
\colorlet{c2}{blue!10}
\colorlet{drawColor}{black!50}
\colorlet{commentColor}{green!70!black!90}

\tikzstyle{oval}=[ellipse, align=center, drop shadow, draw=drawColor, fill=white]
\tikzstyle{rect}=[rectangle, rounded corners=2pt, align=center, drop
shadow, draw=drawColor, fill=white]
\tikzstyle{comment}=[text=commentColor,font=\itshape]
\tikzstyle{textLab}=[]
\tikzstyle{arrow}=[->, very thick, >=stealth', draw=black!80]
\tikzstyle{darrow}=[->, dash pattern=on 3pt off2pt, very thick, >=stealth', draw=black!80]
\tikzstyle{fStartEnd}=[ellipse, align=center, drop shadow, draw=drawColor, fill=white]
\tikzstyle{fInput}=[trapezium, trapezium left angle=70, trapezium right angle=110,
align=center, drop shadow, draw=drawColor, fill=white]
\tikzstyle{fProcess}=[rectangle, align=center, drop shadow, draw=drawColor, fill=white]
\tikzstyle{fSelection}=[diamond, shape aspect=3, align=center, drop
shadow, draw=drawColor, fill=white]
\tikzstyle{fOutput}=[tape, tape bend top=none, align=center, drop shadow, draw=drawColor, fill=white]
\tikzstyle{mem}=[rectangle, align=center, draw=drawColor, fill=white]
\tikzstyle{clo}=[cloud, aspect=2, align=center, drop shadow, draw=drawColor, fill=white]

\lstdefinestyle{customc}{
   language=C,
   % basicstyle=\small\ttfamily\bfseries,
   basicstyle=\ttfamily,
   keywordstyle=\color{blue}\ttfamily,
   stringstyle=\color{red}\ttfamily,
   commentstyle=\color{green}\ttfamily,
   morecomment=[l][\color{magenta}]{\#},
   % breaklines=false,
    breaklines=true, breakatwhitespace=false,
   frameround=fttt,
   frame=trBL,
   backgroundcolor=\color{yellow!20},
   numbers=left,
   stepnumber=1,    
   firstnumber=1,
   numberfirstline=true,
   numberstyle=\tiny\color{black!50},
   xleftmargin=2em,
   framexleftmargin=1.5em
   % linewidth=8cm,
}

\lstnewenvironment{cblock}[1][]
{
  \lstset{
    style=customc,
    #1
  }
}{}

\newcommand{\cfile}[2][]{
  \lstinputlisting[style=customc, #1]{#2}
}

\definecolor{links}{HTML}{2A1B81}
\hypersetup{colorlinks,linkcolor=links,urlcolor=links}

\definecolor{links}{HTML}{2A1B81}
\hypersetup{colorlinks,linkcolor=,urlcolor=links}


\mode<presentation>{
  %-------------------------1
  \usetheme{Boadilla}
  \usecolortheme{beaver}
  %-------------------------1
  %-------------------------2
  %\usetheme{Goettingen}
  %\usecolortheme{sidebartab}
  %-------------------------2
  %\useoutertheme[right]{sidebar}
  %\usefonttheme{default}
  \setbeamercovered{transparent}
  %\setbeameroption{show notes on second screen=right}
  \setbeamertemplate{navigation symbols}{}
  \setbeamertemplate{footline}{}

  \bibliographystyle{abbrv}  
  %\renewcommand\bibfont{\scriptsize}
  \setbeamertemplate{bibliography item}{\textbullet}
  \setbeamertemplate{itemize item}{\checkmark}
  \setbeamertemplate{itemize subitem}{-}
  \setbeamertemplate{enumerate items}[default]
  \setbeamertemplate{sections/subsections in toc}[square]
}

\subtitle{Logical Computational Thinking}
\institute[Tecnológico de Monterrey]{
  \includegraphics[width=5cm]{img/logoTEC.jpg}\\[5mm]
  \includegraphics[width=1cm]{img/logoLEO.jpg}
  Scuola Leonardo Da Vinci (Firenze)
}

\author[Stefano Martina]{
  %\\[0.2cm]
  \textbf{Stefano MARTINA}\\
  {\small stefano.martina@gmail.com}
}

\titlegraphic{\tiny
  \href{http://creativecommons.org/licenses/by-sa/4.0/}{\includegraphics[width=1cm]{img/logoCC.png}}
  This work is licensed under a
  \href{http://creativecommons.org/licenses/by-sa/4.0/}{Creative
    Commons Attribution-ShareAlike 4.0 International License}.}


\title[Lesson 3]{\textbf{Lesson 3 - Compiling a \C\ program}}
\date[10/9/15]{\flushright 10 September 2015}

\begin{document}

\begin{frame}[plain]
  \titlepage
\end{frame}

\begin{frame}
  \frametitle{\texttt{gcc} compiler, GNU/Linux}
  Installed and working by default on every default.
\end{frame}

\begin{frame}
  \frametitle{\texttt{gcc} compiler, Mac}
  \footnotesize
  \begin{block}{Install}
    \begin{enumerate}
    \item Go to \url{https://developer.apple.com/downloads/};
    \item provide authentication, and if needed accept the agreement;
    \item search \alert{command line tools} on the left box;
    \item double click on \alert{Command line tools (...) for Xcode
        6.4};
    \item download \alert{dmg} file (click on the filename on the right);
    \item double click on the \alert{dmg} and then double
      click on the \alert{pkg} inside it;
    \item follow the installation instructions.
    \end{enumerate}
  \end{block}
  \begin{block}{First use}
    \begin{enumerate}
    \item Open a terminal;
    \item type \alert{\texttt{sudo gcc -v}};
    \item provide authentication;
    \item read license, with \alert{space}, and agree, typing
      \alert{\texttt{agree}} at the end (attention to not press to
      much spaces).
    \end{enumerate}
  \end{block}
\end{frame}

\begin{frame}
  \frametitle{\texttt{gcc} compiler, Windows}
  \begin{block}{Install}
    \begin{enumerate}
    \item Go to \url{https://cygwin.com/install.html};
    \item download and run on the corresponding \alert{exe} (32 or 64 bits);
    \item select \alert{Install from internet} and all the default options;
    \item when prompted, select a near server like \alert{\texttt{http://bo.mirror.gar.it}};
    \item when the packages list appear click on the + left to \alert{texttt{devel}} and then on \alert{\texttt{skip}} on the left of \alert{\texttt{gcc-core}} package (the word \texttt{skip} will change to a version number);
    \item click \alert{\texttt{next}} and accept all the question, when finished installing select \alert{\texttt{finish}}.
    \end{enumerate}
  \end{block}
  \begin{block}{First use}
    \begin{enumerate}
    \item Open Cygwin from programs or desktop;
    \item type \alert{\texttt{cd /cygdrive/c/Users/[YourName]/Documents}} and you will go on your documents directory;
    \item go to your preferred working folder where you have sources files (\alert{\texttt{ls}} for listing files and directories, \alert{\texttt{cd}} for changing directory, \alert{\texttt{tab}} autocomplete path);
    \item you can use the \alert{\texttt{gcc}} compiler from here.
    \end{enumerate}
  \end{block}
  \begin{block}{Optional}
    Programs compiled with \alert{\texttt{gcc}} inside Cygwin can be executed only inside Cygwin.
    If you want to compile native Windows programs select also \alert{\texttt{mingw-gcc-core}} package during cygwin installation (you can re-run a cygwin installation and add packages), and then compile sources using \alert{\texttt{i686-pc-mingw32-gcc}} instead of \alert{\texttt{gcc}}.
  \end{block}
\end{frame}
\end{document}
